\section{群}
\label{sec:lie-groups}

重新排列三个物品的位置、把平面整体平移一段距离、把平面绕原点转一个角度，这三类操作作用的对象各不相同，却遵循同样的规则：两个操作可以接连施加，得到的仍是同类操作；“什么都不做”也算一个操作；每个操作都可以撤销。本节先把这些规则从具体对象中抽出来，建立群、子群与同态。这里只固定变换之间怎样复合，同一个群怎样作用在不同对象上留给第~\ref{sec:lie-actions}节，这样就不会把一个群只理解成某一套坐标矩阵。

\subsection{群与子群}
\label{sec:lie-group-subgroup}

本小节先从置换的复合读出群公理，再看几类基本例子和反例，最后研究限制允许变换得到的子群。

\paragraph{从变换的复合到群公理}
先看集合$\{1,2,3\}$上的置换，即从这个集合到自身的双射。置换$\sigma$交换$1$与$2$、保持$3$不动，置换$\tau$交换$2$与$3$、保持$1$不动。复合$\sigma\circ\tau$表示先做$\tau$、再做$\sigma$：
\[
(\sigma\circ\tau)(1)=\sigma(1)=2,\qquad
(\sigma\circ\tau)(2)=\sigma(3)=3,\qquad
(\sigma\circ\tau)(3)=\sigma(2)=1.
\]
反过来先做$\sigma$再做$\tau$，得到$1\mapsto3$、$2\mapsto1$、$3\mapsto2$。两个结果不同，所以复合的次序必须写清。

这个例子已经包含群的全部要素。两个置换的复合仍是置换；恒等置换$\operatorname{id}$不改变任何元素；每个置换都有逆置换，把它的效果撤销；映射的复合满足结合律$(\alpha\circ\beta)\circ\gamma=\alpha\circ(\beta\circ\gamma)$，因为两边在每个元素上都等于$\alpha(\beta(\gamma(i)))$。平面上的平移$\bm x\mapsto\bm x+\bm a$和绕原点的旋转同样满足这四条，只是平移的复合对应向量相加，旋转的复合对应转角相加。

\begin{definition}[群]
\label{def:lie-group}
设$G$是集合，$G\times G\to G$，$(g,h)\mapsto gh$是$G$上的二元运算。若
\begin{enumerate}
 \item 对所有$g,h,k\in G$，$(gh)k=g(hk)$；
 \item 存在$e\in G$，使对所有$g\in G$，$eg=ge=g$；
 \item 对每个$g\in G$，存在$g^{-1}\in G$，使$gg^{-1}=g^{-1}g=e$，
\end{enumerate}
就称$G$连同这个运算构成\term{群}（group），$e$称为单位元，$g^{-1}$称为$g$的逆元。若还对所有$g,h$有$gh=hg$，就称$G$为\term{交换群}（commutative group）。
\end{definition}

运算“封闭”已经写在$G\times G\to G$之中：乘积必须仍在$G$里。乘法记号$gh$只是约定，它可以是映射的复合、数的加法或矩阵乘法。对由变换组成的群，本章总按复合约定$gh=g\circ h$，即先做$h$、再做$g$。

\paragraph{几个基本例子}
$\{1,\ldots,n\}$上全部置换在复合下构成群，称为\term{对称群}（symmetric group），记为$S_n$，它有$n!$个元素。上面的计算说明$S_3$不是交换群。$\mathbb R^d$在向量加法下构成交换群，单位元是$\bm0$，逆元是$-\bm a$；它记录平面或空间中的全部平移。

平面旋转由矩阵
\begin{equation}
 \bm R(\theta)=\begin{bmatrix}\cos\theta&-\sin\theta\\\sin\theta&\cos\theta\end{bmatrix}
 \label{eq:lie-planar-rotation}
\end{equation}
表示。由和角公式，$\bm R(\theta_1)\bm R(\theta_2)=\bm R(\theta_1+\theta_2)$，所以平面旋转的复合与次序无关，单位元是$\bm R(0)=\bm I$，逆元是$\bm R(-\theta)=\bm R(\theta)^\top$。

更一般的矩阵例子如下。全体$n$阶可逆实矩阵在矩阵乘法下构成\term{一般线性群}（general linear group），记为$\mathrm{GL}(n,\mathbb R)$。满足$\bm Q^\top\bm Q=\bm I$的正交矩阵构成\term{正交群}（orthogonal group）$\mathrm O(n)$：若$\bm Q_1,\bm Q_2$正交，则$(\bm Q_1\bm Q_2)^\top\bm Q_1\bm Q_2=\bm Q_2^\top\bm Q_2=\bm I$，且$\bm Q^{-1}=\bm Q^\top$仍正交。正交矩阵的行列式为$\pm1$，其中行列式为$1$的构成\term{特殊正交群}（special orthogonal group）$\mathrm{SO}(n)$。当$n=2$时，$\mathrm{SO}(2)$恰由式~\eqref{eq:lie-planar-rotation}中的矩阵组成：第一列$(a,b)^\top$是单位向量，第二列与它正交且使行列式为$1$，只能是$(-b,a)^\top$，再记$a=\cos\theta$、$b=\sin\theta$即可。

\paragraph{不构成群的情形}
公理中的每一条都会实际排除一些对象。全体$n$阶实矩阵在乘法下不是群，因为奇异矩阵没有逆；把一个向量投影到某条直线上的操作不可撤销，同样不能放进群里。整数在减法下不是群，因为$(1-2)-3\neq1-(2-3)$，结合律失败。若一组图像增强操作接连施加后不再属于这一组，运算不封闭，也不能直接称为群。

单位元和逆元一旦存在就是唯一的，因此记号$e$与$g^{-1}$没有歧义。

\begin{lemma}[单位元与逆元的唯一性]
\label{lem:lie-identity-inverse-unique}
群$G$中单位元唯一；每个$g\in G$的逆元唯一。并且对任意$g,h\in G$，$(gh)^{-1}=h^{-1}g^{-1}$。
\end{lemma}

\begin{proof}
若$e,e'$都是单位元，则$e=ee'=e'$。若$h,h'$都是$g$的逆元，则$h=he=h(gh')=(hg)h'=eh'=h'$。最后，由结合律，$(gh)(h^{-1}g^{-1})=g(hh^{-1})g^{-1}=gg^{-1}=e$，另一侧同理，所以$h^{-1}g^{-1}$是$gh$的逆元。
\end{proof}

最后一式的次序与直觉一致：先穿袜子再穿鞋，撤销时要先脱鞋再脱袜子。对矩阵而言，这就是$(\bm A\bm B)^{-1}=\bm B^{-1}\bm A^{-1}$。

\paragraph{子群}
许多问题只允许一部分变换。例如只允许转$90^\circ$的整数倍，或只允许沿水平方向平移。这些变换自身仍然能复合和撤销，于是在原来的群中又形成一个群。

\begin{definition}[子群]
\label{def:lie-subgroup}
设$G$是群，$H\subseteq G$。若$e\in H$，并且对任意$g,h\in H$都有$gh\in H$与$g^{-1}\in H$，就称$H$是$G$的\term{子群}（subgroup）。
\end{definition}

子群继承了原来的运算，结合律自动成立，所以它本身是群。验证子群时，常把三个条件合成一条。

\begin{lemma}[子群判别]
\label{lem:lie-subgroup-criterion}
群$G$的非空子集$H$是子群，当且仅当对任意$g,h\in H$都有$gh^{-1}\in H$。
\end{lemma}

\begin{proof}
必要性由定义直接得到。反过来，取$g\in H$，令$h=g$得$e=gg^{-1}\in H$；再取第一个元素为$e$，得$h^{-1}=eh^{-1}\in H$；最后对$g,h\in H$，由$h^{-1}\in H$得$gh=g(h^{-1})^{-1}\in H$。
\end{proof}

转$90^\circ$整数倍的旋转$\{\bm R(k\pi/2):k=0,1,2,3\}$是$\mathrm{SO}(2)$的子群，它只有四个元素。水平平移$\{(a,0)^\top:a\in\mathbb R\}$是$\mathbb R^2$的子群。矩阵群之间有包含链$\mathrm{SO}(n)\subseteq\mathrm O(n)\subseteq\mathrm{GL}(n,\mathbb R)$，每一个都是后一个的子群。

\paragraph{子集与子群的区别}
子群是“允许的变换”组成的子集，并且对群运算封闭；它不同于被变换对象的一个子集，也不同于群的任意子集。在平移群$\mathbb R^2$中，单位圆$\{\bm x:\lVert\bm x\rVert_2=1\}$是一个子集，却不是子群：$(1,0)^\top+(1,0)^\top=(2,0)^\top$不在圆上，零向量也不在圆上。

同样，$\mathrm O(n)$中行列式为$-1$的矩阵（反射类变换）组成的子集不是子群：它不含$\bm I$，两个反射的乘积行列式为$1$，又回到了旋转。可见“满足某个条件的变换”能否形成子群，要逐条核对单位元、乘积和逆元，不能只看条件写起来是否自然。

\subsection{群同态}
\label{sec:lie-homomorphism}

同一套变换规则可以用不同方式记录。平面旋转可以用转角$\theta$记录，也可以用矩阵$\bm R(\theta)$记录；三个物品的置换可以写成映射表，也可以写成置换矩阵。要说明两种记录描述的是同一个复合结构，需要一个保持复合的映射。

\begin{definition}[群同态与同构]
\label{def:lie-homomorphism}
设$G,H$是群。映射$\varphi:G\to H$若对所有$g_1,g_2\in G$满足
\begin{equation}
 \varphi(g_1g_2)=\varphi(g_1)\varphi(g_2),
 \label{eq:lie-homomorphism}
\end{equation}
就称为\term{群同态}（group homomorphism）。双射的群同态称为\term{群同构}（group isomorphism）。同态的\term{核}（kernel）是$\ker\varphi=\{g\in G:\varphi(g)=e_H\}$，像是$\varphi(G)=\{\varphi(g):g\in G\}$。
\end{definition}

式~\eqref{eq:lie-homomorphism}左边先在$G$中复合再映过去，右边先映过去再在$H$中复合，二者相等。由$\varphi(e_G)=\varphi(e_Ge_G)=\varphi(e_G)\varphi(e_G)$，两边左乘$\varphi(e_G)^{-1}$得$\varphi(e_G)=e_H$；再由$\varphi(g)\varphi(g^{-1})=\varphi(e_G)=e_H$得$\varphi(g^{-1})=\varphi(g)^{-1}$。所以同态自动保持单位元与逆元。

\paragraph{同态的例子}
由式~\eqref{eq:lie-planar-rotation}定义的$\theta\mapsto\bm R(\theta)$是从加法群$\mathbb R$到$\mathrm{SO}(2)$的同态，因为$\bm R(\theta_1+\theta_2)=\bm R(\theta_1)\bm R(\theta_2)$。它是满射，核为$\{2\pi k:k\in\mathbb Z\}$：转角相差$2\pi$的整数倍，得到同一个旋转。

行列式给出同态$\det:\mathrm{GL}(n,\mathbb R)\to\mathbb R\setminus\{0\}$，后者取数的乘法，因为$\det(\bm A\bm B)=\det\bm A\det\bm B$。它的核是行列式为$1$的矩阵全体，记为$\mathrm{SL}(n,\mathbb R)$。

对$\pi\in S_n$，令置换矩阵$\bm P_\pi$的第$j$列为$\bm e_{\pi(j)}$，即$\bm P_\pi\bm e_j=\bm e_{\pi(j)}$。于是
\[
\bm P_{\pi_1}\bm P_{\pi_2}\bm e_j=\bm P_{\pi_1}\bm e_{\pi_2(j)}=\bm e_{\pi_1(\pi_2(j))}=\bm P_{\pi_1\pi_2}\bm e_j,
\]
所以$\pi\mapsto\bm P_\pi$是从$S_n$到$\mathrm{GL}(n,\mathbb R)$的同态，像落在$\mathrm O(n)$中。不同置换给出不同矩阵，这个同态是单射。

把群的复合规则落实为可逆线性映射，是比较各种表示时最常用的同态，下面单独命名。

\begin{definition}[线性表示]
\label{def:lie-linear-representation}
设$G$是群，$V$是有限维实向量空间，$\mathrm{GL}(V)$是$V$上全部可逆线性映射在复合下构成的群。群同态$\rho:G\to\mathrm{GL}(V)$称为$G$在$V$上的\term{线性表示}（linear representation）。
\end{definition}

同一个群可以有许多线性表示。$\mathrm{SO}(2)$可以按$\rho(\bm R)=\bm R$作用在$\mathbb R^2$上，也可以按$\rho(\bm R)=1$作用在$\mathbb R$上，后者把每个旋转都表示为“不做任何事”。第~\ref{sec:lie-invariant-equivariant}节会用这两种表示区分不变与等变。

\paragraph{核与像保留什么}
核记录同态把哪些元素“看成单位元”，因而丢失了哪些区分；像记录目标群中有哪些元素被实际表示出来。下面的结论说明，核只含单位元时同态不丢失任何区分。

\begin{theorem}[核与单射]
\label{thm:lie-kernel-injective}
设$\varphi:G\to H$是群同态。则$\ker\varphi$是$G$的子群，$\varphi(G)$是$H$的子群，并且$\varphi(g_1)=\varphi(g_2)$当且仅当$g_1^{-1}g_2\in\ker\varphi$。特别地，$\varphi$是单射当且仅当$\ker\varphi=\{e_G\}$。
\end{theorem}

\begin{proof}
若$g,h\in\ker\varphi$，则$\varphi(gh^{-1})=\varphi(g)\varphi(h)^{-1}=e_H$，由引理~\ref{lem:lie-subgroup-criterion}知核是子群；对像中的$\varphi(g),\varphi(h)$，有$\varphi(g)\varphi(h)^{-1}=\varphi(gh^{-1})$仍在像中，所以像也是子群。其次，$\varphi(g_1)=\varphi(g_2)$等价于$e_H=\varphi(g_1)^{-1}\varphi(g_2)=\varphi(g_1^{-1}g_2)$，即$g_1^{-1}g_2\in\ker\varphi$。最后，若核只含$e_G$，则$\varphi(g_1)=\varphi(g_2)$推出$g_1^{-1}g_2=e_G$，即$g_1=g_2$；反之，单射把$e_G$之外的元素都映到$e_H$之外。
\end{proof}

回到转角的例子：$\varphi(\theta)=\bm R(\theta)$的核是$2\pi\mathbb Z$，所以$\bm R(\theta_1)=\bm R(\theta_2)$当且仅当$\theta_2-\theta_1\in2\pi\mathbb Z$。转角能完整记录旋转的复合，却在每个旋转上多出可数个等价的记录；第~\ref{sec:lie-rotations}节的轴角与四元数会遇到同类的冗余。

\paragraph{结构对应与集合对应}
两个群元素个数相同，只说明它们之间有集合层面的双射，不说明复合规则相同。$S_3$有$6$个元素，$\{0,1,\ldots,5\}$在模$6$加法下也构成$6$个元素的群，二者之间当然有双射。但模$6$加法群是交换群，$S_3$不是；若$\varphi$是同构，则$\varphi(\sigma)\varphi(\tau)=\varphi(\tau)\varphi(\sigma)$，由单射性得$\sigma\tau=\tau\sigma$，与本节开头的计算矛盾。所以这两个群不同构。

另一方面，平面旋转$\bm R(\theta)$与单位圆上的点$(\cos\theta,\sin\theta)^\top$一一对应。只有给单位圆指定运算，例如把点看成单位复数$e^{\mathrm i\theta}$并取复数乘法，这个对应才成为群同构。判断两种记录是否表达同一个变换结构，要检查的是复合规则是否被保持，而不只是元素能否配对。
